% ID: SEL-D0019-OA-E03
% ESTADO: cotejada. Segunda pasada visual completa y borrador compilado revisados.
% Fuente: Selectividad/matematicas_aplicadas/reserva a/Examen-Matemáticas Aplicadas a las Ciencias Sociales.pdf
% SHA-256: 6fed0384ff6fcbaaa7098ab3f35ea5f881d1280502f050ada3a7f13af528392e
% Página: 1; opción A; ejercicio original 3.

En un centro de enseñanza secundaria el 48\% de los estudiantes son chicos. El 85\% de los chicos del centro y el 82\% de las chicas supera todas las asignaturas. Se elige al azar un estudiante del centro.
\begin{enumerate}[label=\alph*)]
\item \textbf{(1.5 puntos)} \textquestiondown{}Cuál es la probabilidad de que supere todas las asignaturas?
\item \textbf{(1 punto)} Si ha superado todas las asignaturas, \textquestiondown{}cuál es la probabilidad de que sea una chica?
\end{enumerate}
